\section{The codes: small codes, surface and toric codes}\label{sec:codes}

This section and the next go through the code families of \cref{tab:codefamilies} one by one: what the stabilisers are, how a logical state is encoded, how the syndrome is decoded, and (in \cref{sec:ft-clifford,sec:ft-t}) which logical gates the structure supports. The small codes come first because everything about them can be written down completely; the surface code because it is the baseline of every physical estimate in \cref{sec:physical}; colour and quantum LDPC codes because they are what IBM, IonQ and the neutral-atom proposals actually plan to build.

\subsection{Small codes}\label{sec:smallcodes}

\subsubsection{Shor's $\bits{9}{1}{3}$ code}
Shor's code \cite{shor95} is the concatenation of the three-qubit phase-flip code $\ket\pm\mapsto\ket{\pm\pm\pm}$ with the three-qubit bit-flip code $\ket0\mapsto\ket{000}$, $\ket1\mapsto\ket{111}$:
\[
\ket{0_L}=\tfrac1{2\sqrt2}(\ket{000}+\ket{111})^{\otimes3},\qquad
\ket{1_L}=\tfrac1{2\sqrt2}(\ket{000}-\ket{111})^{\otimes3}.
\]
Its eight generators are the six weight-two $Z$-checks $Z_1Z_2,Z_2Z_3,Z_4Z_5,Z_5Z_6,Z_7Z_8,Z_8Z_9$ (bit flips inside each block) and the two weight-six $X$-checks $X_1\cdots X_6$, $X_4\cdots X_9$ (relative phases between blocks). The logical operators are $\bar Z=X_1X_2X_3\ (\equiv X^{\otimes9})$ and $\bar X=Z_1Z_4Z_7\ (\equiv Z^{\otimes9})$---note the swap of letters, a consequence of the outer code being a phase-flip code. The encoder (\cref{fig:small-encoders}, left) is the two concatenation steps written out: two CNOTs spread the data qubit over the three block leaders, Hadamards turn the phase-flip code into the bit-flip basis, and six CNOTs fill each block. Decoding is a lookup (\cref{tab:syn-shor}): the two $Z$-checks of a block locate an $X$ error on one of its three qubits, and the two $X$-checks locate the block carrying a $Z$ error---but not which qubit, because $Z_1$, $Z_2$ and $Z_3$ differ by stabilisers ($Z_1Z_2\in\mathcal S$) and act identically on the code. This is \emph{degeneracy}: three different errors, one syndrome, one correction that works for all three. A $Y$ error is corrected as an $X$ and a $Z$ error simultaneously.

\begin{table}[H]\centering\scriptsize
\begin{tabular}{@{}lccccccccc@{}}\toprule
qubit $j$ & 1 & 2 & 3 & 4 & 5 & 6 & 7 & 8 & 9\\
\midrule
$X_j$ ($Z$-checks, by block) & 10\,00\,00 & 11\,00\,00 & 01\,00\,00 & 00\,10\,00 & 00\,11\,00 & 00\,01\,00 & 00\,00\,10 & 00\,00\,11 & 00\,00\,01\\
$Z_j$ ($X_{1..6},X_{4..9}$) & 10 & 10 & 10 & 11 & 11 & 11 & 01 & 01 & 01\\
\bottomrule\end{tabular}

\caption{Shor-code lookup table (computed, \code{scripts/syndrome\_tables.py}). Each $X$ error has its own syndrome; the three $Z$ errors of a block share one.}
\label{tab:syn-shor}
\end{table}

\subsubsection{Steane's $\bits{7}{1}{3}$ code}
The Steane code \cite{steane96} is the CSS code built from the classical $[7,4,3]$ Hamming code $C$ and its dual: $H_X=H_Z=H_{\rm Ham}$, the $3\times7$ matrix whose $j$-th column is the binary representation of $j$,
\[
H_{\rm Ham}=\begin{pmatrix}0&0&0&1&1&1&1\\0&1&1&0&0&1&1\\1&0&1&0&1&0&1\end{pmatrix},\qquad
\ket{0_L}=\tfrac1{\sqrt8}\sum_{c\in C^\perp}\ket c,\qquad \ket{1_L}=\bar X\ket{0_L},
\]
so $k=7-3-3=1$ and $d=3$ because the Hamming code has distance $3$ (and so does its dual, the simplex code, with distance $4$). Three $X$-type and three $Z$-type weight-four checks; $\bar X=X^{\otimes7}\equiv X_1X_2X_3$ and $\bar Z=Z^{\otimes7}\equiv Z_1Z_4Z_5$. The encoder in \cref{fig:small-encoders} (centre) is the measurement-free version of the standard-form construction: the data qubit's $X$ is first spread to $X_1X_2X_3$ by two CNOTs, then three ancillas in $\ket+$ apply the three $X$-generators by controlled-$X$ fans. Decoding is the Hamming decoder twice (\cref{tab:syn-steane}): the three $Z$-check outcomes read as a binary number are the index of the flipped qubit, and the three $X$-check outcomes the index of the phase-flipped qubit. The Steane code is the smallest 2D colour code (\cref{sec:colour}); its transversal Clifford group is why it appears in every trapped-ion and neutral-atom logical demonstration \cite{ryan-anderson21,bluvstein24}.

\begin{table}[H]\centering\scriptsize
\begin{tabular}{@{}cccccccc@{}}\toprule
qubit $j$ & 1 & 2 & 3 & 4 & 5 & 6 & 7\\
$(s_4,s_2,s_1)$ of $X_j$ (from $Z$-checks) = of $Z_j$ (from $X$-checks) & 001 & 010 & 011 & 100 & 101 & 110 & 111\\
\bottomrule\end{tabular}

\caption{Steane-code lookup table (computed): the syndrome is the binary index of the faulty qubit, for $X$ and $Z$ errors alike.}
\label{tab:syn-steane}
\end{table}

\subsubsection{The perfect $\bits{5}{1}{3}$ code}
The five-qubit code \cite{laflamme96,bennett96} is the smallest code correcting an arbitrary single-qubit error: its $16$ syndromes exactly cover the identity and the $15$ single-qubit Paulis (\cref{tab:syn-513}), which is the quantum Hamming bound with equality. Its stabiliser group is generated by $XZZXI$ and its cyclic shifts $IXZZX$, $XIXZZ$, $ZXIXZ$ (the fifth shift is their product), and $\bar X=X^{\otimes5}$, $\bar Z=Z^{\otimes5}$. It is not CSS, so it has no transversal CNOT and its encoder needs controlled-$Y$ and $S$ gates; \cref{fig:small-encoders} (right) is the standard-form circuit of \cref{alg:encoder}, verified with Stim to map $X_{\rm data}\mapsto+XXXXX$ and $Z_{\rm data}\mapsto+ZZZZZ$ (the leading $Z$ on the data qubit fixes the sign of $\bar X$; two controlled-$Z$ gates from the data qubit to ancillas in $\ket0$ that the generic algorithm emits are omitted because they act trivially).

\begin{table}[H]\centering\scriptsize
\begin{tabular}{@{}lccccccccccccccc@{}}\toprule
$s_1s_2s_3s_4$ & 0001 & 0010 & 0011 & 0100 & 0101 & 0110 & 0111 & 1000 & 1001 & 1010 & 1011 & 1100 & 1101 & 1110 & 1111\\
error & $X_{1}$ & $Z_{3}$ & $X_{5}$ & $Z_{5}$ & $Z_{2}$ & $X_{4}$ & $Y_{5}$ & $X_{2}$ & $Z_{4}$ & $Z_{1}$ & $Y_{1}$ & $X_{3}$ & $Y_{2}$ & $Y_{3}$ & $Y_{4}$\\
\bottomrule\end{tabular}

\caption{$\bits{5}{1}{3}$ lookup table (computed) for generators $S_1=XZZXI$, $S_2=IXZZX$, $S_3=XIXZZ$, $S_4=ZXIXZ$: every one of the $15$ single-qubit errors has a distinct non-zero syndrome.}
\label{tab:syn-513}
\end{table}

\begin{figure}[H]
\centering
\resizebox{\textwidth}{!}{%
\begin{quantikz}[column sep=5pt,row sep=3pt,font=\footnotesize]
\lstick{$\ket\psi$} & \ctrl{3} & \ctrl{6} & \gate{H} & \ctrl{1} & \ctrl{2} & \qw \\
\lstick{$\ket0$}    & \qw & \qw & \qw & \targ{} & \qw & \qw \\
\lstick{$\ket0$}    & \qw & \qw & \qw & \qw & \targ{} & \qw \\
\lstick{$\ket0$}    & \targ{} & \qw & \gate{H} & \ctrl{1} & \ctrl{2} & \qw \\
\lstick{$\ket0$}    & \qw & \qw & \qw & \targ{} & \qw & \qw \\
\lstick{$\ket0$}    & \qw & \qw & \qw & \qw & \targ{} & \qw \\
\lstick{$\ket0$}    & \qw & \targ{} & \gate{H} & \ctrl{1} & \ctrl{2} & \qw \\
\lstick{$\ket0$}    & \qw & \qw & \qw & \targ{} & \qw & \qw \\
\lstick{$\ket0$}    & \qw & \qw & \qw & \qw & \targ{} & \qw
\end{quantikz}%
\hspace{8mm}
\begin{quantikz}[column sep=5pt,row sep=3pt,font=\footnotesize]
\lstick{$\ket\psi$} & \ctrl{1} & \ctrl{2} & \qw & \targ{} & \targ{} & \qw & \qw \\
\lstick{$\ket0$}    & \targ{} & \qw & \qw & \targ{} & \qw & \targ{} & \qw \\
\lstick{$\ket0$}    & \qw & \targ{} & \qw & \qw & \targ{} & \targ{} & \qw \\
\lstick{$\ket0$}    & \qw & \qw & \gate{H} & \ctrl{-3} & \qw & \qw & \qw \\
\lstick{$\ket0$}    & \qw & \qw & \gate{H} & \qw & \ctrl{-4} & \qw & \qw \\
\lstick{$\ket0$}    & \qw & \qw & \gate{H} & \qw & \qw & \ctrl{-4} & \qw \\
\lstick{$\ket0$}    & \qw & \qw & \qw & \targ{}\vqw{-3} & \targ{}\vqw{-2} & \targ{}\vqw{-1} & \qw
\end{quantikz}%
\hspace{8mm}
\begin{quantikz}[column sep=4pt,row sep=3pt,font=\footnotesize]
\lstick{$a_1=\ket0$} & \qw & \gate{H} & \ctrl{1} & \ctrl{3} & \ctrl{4} & \gate{S} & \qw & \qw & \qw & \qw & \qw & \control{} & \qw & \qw & \qw & \control{} & \qw & \qw & \qw \\
\lstick{$a_2=\ket0$} & \qw & \qw & \control{} & \qw & \qw & \qw & \gate{H} & \ctrl{1} & \ctrl{2} & \ctrl{3} & \qw & \qw & \control{} & \qw & \qw & \qw & \qw & \qw \\
\lstick{$a_3=\ket0$} & \qw & \qw & \qw & \qw & \qw & \qw & \qw & \control{} & \qw & \qw & \gate{H} & \ctrl{-2} & \ctrl{-1} & \ctrl{2} & \qw & \qw & \control{} & \qw \\
\lstick{$a_4=\ket0$} & \qw & \qw & \qw & \control{} & \qw & \qw & \qw & \qw & \control{} & \qw & \qw & \qw & \qw & \qw & \gate{H} & \ctrl{-3} & \ctrl{-1} & \ctrl{1} & \gate{S} \\
\lstick{$\ket\psi$}  & \gate{Z} & \qw & \qw & \qw & \gate{Y} & \qw & \qw & \qw & \qw & \targ{} & \qw & \qw & \qw & \targ{} & \qw & \qw & \qw & \gate{Y} & \qw
\end{quantikz}}
\caption{Encoders verified with Stim (\code{scripts/verify\_codes.py}, \code{scripts/standard\_form\_encoder.py}). Left: Shor's code, the two concatenation steps. Centre: Steane's code, data qubit on the top wire; the three $\ket+$ ancillas apply the $X$-type generators $X_{1,2,4,7}$, $X_{1,3,5,7}$, $X_{2,3,6,7}$ (in the figure's wire order the supports are $\{1,2,4,7\},\{1,3,5,7\},\{2,3,6,7\}$). Right: the $\bits{5}{1}{3}$ standard-form encoder; the vertical lines ending in dots are controlled-$Z$ gates, the boxed $Y$ is a controlled-$Y$, and each $S$ supplies the phase for a generator with $Y$ on its pivot ($YZIZY$ and $ZIZYY$, which together with $IXZZX$ and $ZZXIX$ generate the same group as the cyclic generators).}
\label{fig:small-encoders}
\end{figure}

\subsubsection{The $\bits{15}{1}{3}$ Reed--Muller code and the $\bits{4}{2}{2}$ detection code}
Index the $15$ qubits by the non-zero $4$-bit strings $v\in\F_2^4\setminus\{0\}$. The $X$-stabilisers are the four weight-$8$ operators $X$ on $\{v:v_b=1\}$, $b=1..4$; the $Z$-stabilisers are those four plus the six weight-$4$ operators $Z$ on $\{v:v_b=v_{b'}=1\}$; $\bar X=X^{\otimes15}$, $\bar Z=Z^{\otimes15}\equiv$ a weight-three $Z$ (so $d=3$, while the $X$-distance is $7$). The CSS structure was checked and a nine-instruction standard-form encoder verified (\cref{app:scripts}). What makes this code important is that every $X$-stabiliser has weight $\equiv0\pmod8$ (the classical code is \emph{triply even}) and $\bar X$ has weight $15\equiv7\pmod 8$: writing $\ket{0_L}\propto\sum_{s\in\mathcal S_X}\ket s$ and $\ket{1_L}\propto\sum_s\ket{s\oplus1^{15}}$, the transversal gate $T^{\otimes15}$ multiplies a basis state $\ket x$ by $e^{i\pi|x|/4}$, which is $1$ on every term of $\ket{0_L}$ and $e^{-i\pi/4}$ on every term of $\ket{1_L}$: $T^{\otimes15}=\bar T^\dagger$. This is the engine of magic-state distillation (\cref{sec:distillation}) and of code switching (\cref{sec:codeswitch}); the code is also the smallest 3D colour code.

The $\bits{4}{2}{2}$ code, $\mathcal S=\langle XXXX,ZZZZ\rangle$, $\bar X_1=X_1X_2$, $\bar Z_1=Z_1Z_3$, $\bar X_2=X_1X_3$, $\bar Z_2=Z_1Z_2$, detects any single error (every weight-one Pauli anticommutes with one of the two checks) but cannot locate it; it is the smallest member of the $\bits{2m}{2m-2}{2}$ ``iceberg'' family used in trapped-ion and neutral-atom demonstrations where post-selection on a trivial syndrome is affordable \cite{bluvstein24}.

\subsection{The toric and surface codes}\label{sec:surface}

\subsubsection{Kitaev's toric code}
Place qubits on the $2L^2$ edges of an $L\times L$ square lattice with periodic boundary conditions \cite{kitaev03,dennis02}. For each vertex $v$ and each face (plaquette) $p$ define
\[
A_v=\prod_{e\ni v}X_e\quad(\text{weight }4),\qquad B_p=\prod_{e\in\partial p}Z_e\quad(\text{weight }4).
\]
A vertex and a plaquette share either $0$ or $2$ edges, so all generators commute. Since $\prod_vA_v=\prod_pB_p=I$, there are $2(L^2-1)$ independent generators and $k=2L^2-2(L^2-1)=2$ logical qubits. A product of $Z$'s along a closed loop of edges commutes with every $A_v$; if the loop is contractible it is a product of $B_p$'s (a stabiliser), if it winds around the torus it is a logical operator. The two non-contractible cycles give $\bar Z_1,\bar Z_2$; $\bar X_1,\bar X_2$ are $X$-strings along non-contractible cycles of the \emph{dual} lattice (crossing edges). The shortest non-contractible loop has $L$ edges, so $d=L$: the toric code is $\bits{2L^2}{2}{L}$. Its decoding problem is the prototype for everything below: a $Z$-error on a set of edges is a chain; $A_v$ is violated exactly at the \emph{endpoints} of the chain (an interior vertex meets the chain twice); so the syndrome is a set of defects, always even in number, and the decoder's job is to pair them by a correction chain such that error$+$correction is a contractible loop---a stabiliser---rather than a non-contractible one.

\subsubsection{The planar (rotated) surface code}
Cut the torus open and the two logical qubits become one, held by a patch with two kinds of boundary \cite{bravyi-kitaev98,fowler12}. In the \emph{rotated} layout of \cref{fig:surface-lattice}, $d^2$ data qubits sit on a $d\times d$ grid, and the $d^2-1$ checks are the faces of the grid, coloured in a checkerboard: $X$-type faces measure $X$ on their (four, or two at the boundary) corners, $Z$-type faces measure $Z$. Two opposite boundaries carry weight-two $X$ checks, the other two carry weight-two $Z$ checks, and a logical operator is a string that runs between the two boundaries of its own type: $\bar Z$ is a row of $Z$'s between the $Z$ boundaries (each bulk $X$ face it crosses is touched at two corners, so it commutes with everything), $\bar X$ a column of $X$'s between the $X$ boundaries. The shortest such string has $d$ qubits; deforming a string by multiplying with a face operator gives an equivalent representative, so each logical operator has many representatives, and this high degeneracy is also why the code decodes so well. One patch is $\bits{d^2}{1}{d}$ on $2d^2-1$ physical qubits once the $d^2-1$ measurement ancillas are counted.

\begin{figure}[H]
\centering
\begin{tikzpicture}[scale=1.05]
\def\d{5}
% bulk faces
\foreach \i in {0,...,3}{\foreach \j in {0,...,3}{
  \pgfmathparse{int(mod(\i+\j,2))}
  \ifnum\pgfmathresult=0
    \fill[xface] (\j,-\i) rectangle (\j+1,-\i-1);
  \else
    \fill[zface] (\j,-\i) rectangle (\j+1,-\i-1);
  \fi}}
% boundary half-discs: top X at columns 1-2 and 3-4 ; bottom X at 0-1 and 2-3; left Z rows 0-1, 2-3 ; right Z rows 1-2, 3-4
\foreach \j in {1,3}{\fill[xface] (\j+1,0) arc (0:180:0.5);}
\foreach \j in {0,2}{\fill[xface] (\j,-4) arc (180:360:0.5);}
\foreach \i in {0,2}{\fill[zface] (0,-\i) arc (90:270:0.5);}
\foreach \i in {1,3}{\fill[zface] (4,-\i-1) arc (270:450:0.5);}
% grid lines
\draw[black!40,thin] (0,0) grid[step=1] (4,-4);
% logical operators
\draw[line width=2.2pt,cZ!80!black,opacity=0.85] (0,-2) -- (4,-2);
\node[font=\scriptsize,text=cZ!80!black] at (4.9,-2) {$\bar Z$};
\draw[line width=2.2pt,cX!80!black,opacity=0.85] (1,0) -- (1,-4);
\node[font=\scriptsize,text=cX!80!black] at (1,0.45) {$\bar X$};
% data qubits
\foreach \i in {0,...,4}{\foreach \j in {0,...,4}{\node[dq] at (\j,-\i) {};}}
% measure qubits (face centres)
\foreach \i in {0,...,3}{\foreach \j in {0,...,3}{\node[circle,fill=black!70,inner sep=0pt,minimum size=3pt] at (\j+0.5,-\i-0.5) {};}}
\foreach \j in {1,3}{\node[circle,fill=black!70,inner sep=0pt,minimum size=3pt] at (\j+0.5,0.3) {};}
\foreach \j in {0,2}{\node[circle,fill=black!70,inner sep=0pt,minimum size=3pt] at (\j+0.5,-4.3) {};}
\foreach \i in {0,2}{\node[circle,fill=black!70,inner sep=0pt,minimum size=3pt] at (-0.3,-\i-0.5) {};}
\foreach \i in {1,3}{\node[circle,fill=black!70,inner sep=0pt,minimum size=3pt] at (4.3,-\i-0.5) {};}
% legend
\fill[xface] (5.6,-0.2) rectangle (5.95,-0.55); \node[font=\scriptsize,anchor=west] at (6.0,-0.37) {$X$ check ($XXXX$ / $XX$)};
\fill[zface] (5.6,-0.8) rectangle (5.95,-1.15); \node[font=\scriptsize,anchor=west] at (6.0,-0.97) {$Z$ check ($ZZZZ$ / $ZZ$)};
\node[dq] at (5.78,-1.55) {}; \node[font=\scriptsize,anchor=west] at (6.0,-1.55) {data qubit ($d^2=25$)};
\node[circle,fill=black!70,inner sep=0pt,minimum size=3pt] at (5.78,-2.05) {}; \node[font=\scriptsize,anchor=west] at (6.0,-2.05) {measure qubit ($d^2-1=24$)};
\node[font=\scriptsize,anchor=west,text width=5.2cm,align=left] at (5.55,-3.2) {Strings of the same type as a boundary end on it. A $Z$ chain anticommutes only with the $X$ faces at its two ends; the decoder pairs such \emph{defects}, with the boundary as an extra partner.};
% an error chain example
\draw[line width=1.6pt,black,dashed] (1,-3) -- (2,-3);
\node[star,star points=5,star point ratio=2.2,fill=yellow!80!orange,draw=black!60,inner sep=1.2pt] at (0.5,-2.5) {};
\node[star,star points=5,star point ratio=2.2,fill=yellow!80!orange,draw=black!60,inner sep=1.2pt] at (2.5,-2.5) {};
\end{tikzpicture}
\caption{The rotated surface code at $d=5$. The dashed chain is a $Z$ error on two neighbouring qubits of the fourth row; the two stars are the only $X$ faces it anticommutes with (the $X$ face between them is touched twice)---the \emph{defects} the decoder sees. The logical operators drawn are the minimum-weight representatives; multiplying $\bar Z$ by any $Z$ face slides it.}
\label{fig:surface-lattice}
\end{figure}

\paragraph{Encoding.} A logical $\ket{0_L}$ is prepared by initialising all data qubits in $\ket0$ and running one cycle of syndrome measurement: the $Z$ faces are already satisfied, the $X$ faces return random outcomes that are simply recorded as the reference frame (or fixed by $Z$ corrections), and from then on the code is live. $\ket{+_L}$ is the same with $\ket+$ and the $X$ faces. Reading out $\bar Z$ is the time-reversed operation: measure every data qubit in the $Z$ basis, recompute the $Z$ faces from the outcomes to get one last syndrome round (so that the readout itself is error-corrected), and take the parity along any $\bar Z$ string.

\paragraph{Syndrome extraction.} Each face has one measurement qubit. One cycle (\cref{fig:surface-cycle}) is: reset the ancillas; Hadamard the $X$-face ancillas; four layers of CNOTs in which every face talks to its four corners in a fixed order; Hadamard; measure. The order matters because of \emph{hook errors}: an ancilla fault between the second and third CNOT propagates to the last two data qubits of the order, and if that pair were aligned with the logical operator the pair would count as a single fault that advances a chain by two, halving the effective distance. The standard orders---``Z'' shape (NW, NE, SW, SE) for the $X$ faces and ``N'' shape (NW, SW, NE, SE) for the $Z$ faces---make every hook pair perpendicular to the logical string it could have helped, and they also let all faces run their CNOTs simultaneously, because at each layer every data qubit is touched by exactly one face. The whole cycle is $8$ layers; Google's Willow runs it in $1.1\,\mu$s \cite{willow25}, which is the ``$1\,\mu$s cycle'' of every superconducting estimate in \cref{sec:physical}.

\begin{figure}[H]
\centering
\begin{quantikz}[column sep=6pt,row sep=4pt,font=\small]
\lstick{$X$-face ancilla $\ket0$} & \gate{H} & \ctrl{1} & \ctrl{2} & \ctrl{3} & \ctrl{4} & \gate{H} & \meter{} \\
\lstick{NW} & \qw & \targ{} & \qw & \qw & \qw & \qw & \qw \\
\lstick{NE} & \qw & \qw & \targ{} & \qw & \qw & \qw & \qw \\
\lstick{SW} & \qw & \qw & \qw & \targ{} & \qw & \qw & \qw \\
\lstick{SE} & \qw & \qw & \qw & \qw & \targ{} & \qw & \qw
\end{quantikz}
\hspace{6mm}
\begin{quantikz}[column sep=6pt,row sep=4pt,font=\small]
\lstick{$Z$-face ancilla $\ket0$} & \qw & \targ{} & \targ{} & \targ{} & \targ{} & \qw & \meter{} \\
\lstick{NW} & \qw & \ctrl{-1} & \qw & \qw & \qw & \qw & \qw \\
\lstick{SW} & \qw & \qw & \ctrl{-2} & \qw & \qw & \qw & \qw \\
\lstick{NE} & \qw & \qw & \qw & \ctrl{-3} & \qw & \qw & \qw \\
\lstick{SE} & \qw & \qw & \qw & \qw & \ctrl{-4} & \qw & \qw
\end{quantikz}
\caption{One surface-code cycle for an $X$ face (left, order NW--NE--SW--SE) and a $Z$ face (right, order NW--SW--NE--SE). The four CNOT layers of all faces run in parallel; the hook pairs (last two targets) are horizontal for $X$ faces and vertical for $Z$ faces, perpendicular to $\bar X$ and $\bar Z$ respectively in \cref{fig:surface-lattice}.}
\label{fig:surface-cycle}
\end{figure}

\subsubsection{Decoding the surface code}\label{sec:mwpm}
Because measurements are noisy, one round is never enough: the cycle is repeated, and the decoder works on \emph{detectors}, the XOR of consecutive outcomes of the same face \cite{dennis02}. In the resulting three-dimensional space-time picture (\cref{fig:matching}), every elementary fault flips at most two detectors: a data error between rounds flips the two faces beside it (a space-like edge), a measurement error flips the same face in two consecutive rounds (a time-like edge), a CNOT fault or hook flips a diagonal pair (a space-time edge), and a fault next to a boundary flips a single detector (an edge to the boundary node). The \emph{detector error model} is the list of these edges with their probabilities; Stim derives it from the circuit, and \cref{fig:threshold} was produced that way \cite{gidney21stim}. Decoding then means choosing a set of edges that explains the observed defects and is most probable:

\begin{algorithm}[H]
\caption{Minimum-weight perfect matching (MWPM) decoding of a surface-code memory}\label{alg:mwpm}
\begin{algorithmic}[1]
\Require detector error model $G=(V,E)$ with edge probabilities $p_e$ and logical-flip labels; observed defect set $D\subseteq V$ (odd $|D|$ is made even by a boundary node)
\State weight every edge $w_e=\log\frac{1-p_e}{p_e}$ (so that a path's weight is $-\log$ of its probability up to a constant)
\State for every pair of defects, the lightest path between them is the shortest path in $(G,w)$ (Dijkstra); build the complete graph on $D$ with these distances, plus a boundary partner per defect
\State find a \emph{perfect matching} of minimum total weight (Edmonds' blossom algorithm; sparse blossom explores the paths lazily and runs in near-linear time at low $p$ \cite{higgott23pymatching})
\State the correction is the XOR of the matched paths' logical labels: the decoder never touches the qubits, it flips the Pauli frame
\end{algorithmic}
\end{algorithm}

MWPM finds the single most likely error, not the most likely \emph{class} of errors; it ignores degeneracy and the correlations between $X$ and $Z$ decoding caused by $Y$ errors, and each of those costs a little threshold. Fowler showed that matching can be made $O(1)$ average time per round with parallel hardware \cite{fowler13mwpm}; PyMatching's sparse blossom decodes about a million detection events per second per core \cite{higgott23pymatching}; Google's real-time decoder kept up with the $1.1\,\mu$s cycle of a distance-$5$ memory for a million cycles with $63\,\mu$s average latency \cite{willow25}, and \emph{latency}, not throughput, is what sets the reaction time that multiplies every non-Clifford gate in \cref{sec:physical}.

\begin{figure}[H]
\centering
\begin{tikzpicture}[scale=0.9]
% space-time lattice of detectors for one face-type: 5 columns x 4 rounds
\foreach \t in {0,...,3}{\foreach \x in {0,...,4}{\node[circle,draw=black!50,fill=white,inner sep=0pt,minimum size=4pt] (n\x\t) at (\x*1.2,\t*1.0) {};}}
\foreach \t in {0,...,3}{\draw[black!25] (0,\t) -- (4.8,\t);}
\foreach \x in {0,...,4}{\draw[black!25] (\x*1.2,0) -- (\x*1.2,3);}
\draw[black!25,dashed] (-0.8,0) -- (-0.8,3); \node[font=\scriptsize,rotate=90] at (-1.05,1.5) {boundary};
\draw[black!25,dashed] (5.6,0) -- (5.6,3); \node[font=\scriptsize,rotate=90] at (5.85,1.5) {boundary};
% faults: data error at t=1 between x=1 and x=2 -> defects (1,1),(2,1)
\draw[cZ,line width=2pt] (1.2,1) -- (2.4,1);
\node[star,star points=5,star point ratio=2.2,fill=yellow!80!orange,draw=black!60,inner sep=1pt] at (1.2,1) {};
\node[star,star points=5,star point ratio=2.2,fill=yellow!80!orange,draw=black!60,inner sep=1pt] at (2.4,1) {};
% measurement error at x=4 between t=1 and t=2 -> defects (4,1),(4,2)
\draw[cX,line width=2pt] (4.8,1) -- (4.8,2);
\node[star,star points=5,star point ratio=2.2,fill=yellow!80!orange,draw=black!60,inner sep=1pt] at (4.8,1) {};
\node[star,star points=5,star point ratio=2.2,fill=yellow!80!orange,draw=black!60,inner sep=1pt] at (4.8,2) {};
% hook / space-time diagonal at x=0: t=2 -> boundary edge
\draw[black!60,line width=2pt] (0,2.5) -- (-0.8,2.5);
\node[star,star points=5,star point ratio=2.2,fill=yellow!80!orange,draw=black!60,inner sep=1pt] at (0,3) {};
\draw[violet,line width=2pt,dotted] (0,3) -- (0,2.5);
\node[font=\scriptsize,anchor=west] at (6.3,2.8) {$\bullet$ detector (face $f$, round $t$) $=m_{f,t}\oplus m_{f,t-1}$};
\node[font=\scriptsize,anchor=west,text=cZ] at (6.3,2.2) {--- data error: two neighbouring faces, same round};
\node[font=\scriptsize,anchor=west,text=cX] at (6.3,1.6) {--- measurement error: same face, consecutive rounds};
\node[font=\scriptsize,anchor=west,text=black!60] at (6.3,1.0) {--- boundary fault: one defect};
\node[font=\scriptsize,anchor=west,text=violet] at (6.3,0.4) {$\cdots$ matched correction};
\node[font=\scriptsize] at (2.4,-0.5) {space (faces of one type) $\rightarrow$};
\node[font=\scriptsize,rotate=90] at (-1.7,1.5) {rounds $\rightarrow$};
\end{tikzpicture}
\caption{The decoding graph of one check type over four rounds: every fault is an edge, every defect a vertex, and MWPM pairs the defects by lightest edges. A chain of faults that crosses the whole patch (or a matched correction that does) flips the logical operator.}
\label{fig:matching}
\end{figure}

\paragraph{Union-find.}\label{sec:unionfind} Delfosse and Nickerson's decoder \cite{delfosse-nickerson21} trades a little accuracy for almost-linear worst-case time ($O(n\,\alpha(n))$, $\alpha$ the inverse Ackermann function) and a structure that maps to simple hardware:

\begin{algorithm}[H]
\caption{Union-find decoder}\label{alg:uf}
\begin{algorithmic}[1]
\State start one cluster per defect; a cluster is \emph{odd} if it contains an odd number of defects and does not touch the boundary
\While{some cluster is odd}
  \State grow every odd cluster by half an edge in all directions (\emph{syndrome validation})
  \State when two clusters touch, merge them (union-find with path compression keeps this near constant time); a merged cluster's parity is the sum
\EndWhile
\State inside each cluster, build a spanning forest and \emph{peel} it from the leaves: remove a leaf edge, flip it in the correction iff the leaf vertex is an unmatched defect, and toggle the parity of the other endpoint
\end{algorithmic}
\end{algorithm}
Its threshold is a few tenths of a percent below MWPM's under the same noise, which is why it is used where decoding hardware is the constraint and why many real-time decoders start from it.

\paragraph{Learned decoders.} AlphaQubit \cite{alphaqubit24} (Google DeepMind and Google Quantum AI) is a recurrent transformer trained first on simulated data and then fine-tuned on Sycamore syndrome records; it beat correlated matching on the experimental distance-$3$ and $-5$ data by using the soft measurement information and the device's real error correlations, at the price of a latency that is not yet compatible with real-time use.

\subsubsection{The threshold, measured}\label{sec:threshold}
\Cref{fig:threshold} is a circuit-level simulation performed for this document with Stim and PyMatching (\code{surface\_threshold.py}, \cref{app:scripts}): rotated memory-$Z$ patches of distance $3$--$9$, $d$ rounds each, uniform depolarising noise of strength $p$ after every Clifford, before every round on the data, and on every measurement and reset; at least $1{,}000$ logical failures or $12$ million shots per point; the logical error per shot $\epsilon$ converted to an error per round by $\epsilon_{\rm round}=\frac12(1-(1-2\epsilon)^{1/d})$. The curves cross at $p\approx1.1$--$1.2\%$, the circuit-level threshold under this noise model; below it, each $+2$ of distance divides the error by $\Lambda$, which the data put at $\Lambda\approx9.6$, $10.6$ and $14.8$ for the steps $3\to5$, $5\to7$ and $7\to9$ at $p=10^{-3}$, $\approx5$ at $p=2\times10^{-3}$ and $\approx3.2$ at $p=3\times10^{-3}$. Google's Willow reached $\Lambda=2.14\pm0.02$ with a $d=7$ logical error of $0.143\%$ per cycle on $101$ qubits, which in this uniform model corresponds to $p\approx0.45\%$ (real noise is far from uniform: Willow's two-qubit gates are better than that, its measurements worse); the $\Lambda\approx10$ of the $p=10^{-3}$ estimates is a factor four or five in physical error away from that device \cite{willow25}.

\begin{figure}[H]
\centering
\begin{tikzpicture}
\begin{loglogaxis}[width=0.62\textwidth,height=0.5\textwidth,
  xlabel={physical error rate $p$ (circuit-level depolarising)}, ylabel={logical error per round $\epsilon_{\rm round}$},
  xmin=0.0004,xmax=0.02,ymin=1e-8,ymax=0.2, grid=both,grid style={black!10},
  legend style={font=\scriptsize,at={(0.03,0.97)},anchor=north west,cells={anchor=west}},
  tick label style={font=\footnotesize},label style={font=\small},
  xtick={0.0005,0.001,0.002,0.005,0.01,0.02}, xticklabels={$0.05\%$,$0.1\%$,$0.2\%$,$0.5\%$,$1\%$,$2\%$}]
\addplot[color=cZ!70!black,mark=*,mark size=1.4pt,thick] table[x=p,y=eps_round,discard if not={d}{3}] {data/surface_threshold.dat}; \addlegendentry{$d=3$}
\addplot[color=boxgreen,mark=square*,mark size=1.4pt,thick] table[x=p,y=eps_round,discard if not={d}{5}] {data/surface_threshold.dat}; \addlegendentry{$d=5$}
\addplot[color=orange!85!black,mark=triangle*,mark size=1.8pt,thick] table[x=p,y=eps_round,discard if not={d}{7}] {data/surface_threshold.dat}; \addlegendentry{$d=7$}
\addplot[color=boxred,mark=diamond*,mark size=1.8pt,thick] table[x=p,y=eps_round,discard if not={d}{9}] {data/surface_threshold.dat}; \addlegendentry{$d=9$}
\draw[black!50,dashed] (axis cs:0.001,1e-8) -- (axis cs:0.001,0.2);
\node[font=\scriptsize,anchor=south west,rotate=90] at (axis cs:0.001,2e-8) {$p=10^{-3}$ (Model S)};
\draw[black!50,dashed] (axis cs:0.0115,1e-8) -- (axis cs:0.0115,0.2);
\node[font=\scriptsize,anchor=south east,rotate=90] at (axis cs:0.0115,2e-8) {crossing $\approx1.15\%$};
\end{loglogaxis}
\end{tikzpicture}
\hfill
\begin{minipage}[b]{0.34\textwidth}\footnotesize\centering
\begin{tabular}{@{}lccc@{}}\toprule
$p$ & $\Lambda_{3\to5}$ & $\Lambda_{5\to7}$ & $\Lambda_{7\to9}$\\\midrule
$10^{-3}$ & 9.6 & 10.6 & 14.8\\
$2\times10^{-3}$ & 5.3 & 4.8 & 5.0\\
$3\times10^{-3}$ & 3.3 & 3.2 & 3.2\\
$5\times10^{-3}$ & 2.0 & 1.9 & 1.9\\\bottomrule
\end{tabular}\\[4pt]
$\epsilon_{\rm round}(d{=}9,p{=}10^{-3})=1.7\times10^{-7}$ ($18$ failures)\\
extrapolated to $d=25$ at $\Lambda\approx10$: $\approx10^{-15}$\\[6pt]
\end{minipage}
\caption{Surface-code threshold computed for this document (\cref{app:scripts}). Left: logical error per round against physical error for four distances; the crossing is the threshold. Right: the suppression factor $\Lambda$ between successive distances, read off the same data; the extrapolation to $d=25$ uses $\Lambda\approx10$ per step, consistent with \cref{eq:plaw} at $p_{\rm th}\approx1.2\%$. Points with fewer than $200$ failures ($d=7,9$ at $p\le10^{-3}$) carry $\pm10$--$25\%$ statistical error; the $d=9$ point at $p=0.05\%$ had no failure in $12$ million shots and is not drawn.}
\label{fig:threshold}
\end{figure}

\paragraph{Lattice surgery and logical gates.} The surface code has no transversal gate beyond the Paulis, so its logical CNOTs, Hadamards and measurements are made by merging and splitting patches---\emph{lattice surgery}---and its $T$ gates by consuming magic states; both are in \cref{sec:ft-clifford,sec:ft-t}.
